\textbf{1. Kiến thức cần nhớ}
\begin{itemize}
\item \textbf{Quy tắc nhân đơn thức với đa thức}: Muốn nhân một đơn thức với một đa thức, ta nhân đơn thức với từng hạng tử của đa thức rồi cộng các tích với nhau:
80493\boxed{A(B + C) = AB + AC}80493
\item \textbf{Quy tắc nhân đa thức với đa thức}: Muốn nhân một đa thức với một đa thức, ta nhân mỗi hạng tử của đa thức này với từng hạng tử của đa thức kia rồi cộng các tích với nhau:
80493\boxed{(A + B)(C + D) = AC + AD + BC + BD}80493
\item \textbf{Quy tắc lũy thừa}:
80493x^m \cdot x^n = x^{m+n};\quad (x^m)^n = x^{m \cdot n};\quad (xy)^n = x^n y^n.80493
\item \textbf{Quy tắc dấu}:  \cdot (+) = (+)$;\quad  \cdot (-) = (+)$;\quad  \cdot (-) = (-)$.
\end{itemize}

\textbf{2. Phương pháp giải các dạng}

%% DANG 8D1?4-1
\textbf{Dạng 1. Nhân đơn thức với đa thức}
\begin{enumerate}
\item Lấy đơn thức nhân lần lượt với từng hạng tử của đa thức (chú ý nhân cả dấu).
\item Nhân hệ số với hệ số, nhân các lũy thừa cùng cơ số bằng cách cộng số mũ.
\item Thu gọn biểu thức kết quả (nếu có hạng tử đồng dạng).
\end{enumerate}
Lưu ý: Cẩn thận dấu khi đơn thức có dấu trừ, ví dụ: 569X2x(x - 3) = -2x^2 + 6x$.

%% DANG 8D1?4-2
\textbf{Dạng 2. Nhân đa thức với đa thức}
\begin{enumerate}
\item Lấy từng hạng tử của đa thức thứ nhất nhân với toàn bộ đa thức thứ hai.
\item Khai triển thành tổng các tích đơn thức.
\item \textbf{Thu gọn các hạng tử đồng dạng} để được đa thức kết quả dạng tối giản.
\end{enumerate}

%% DANG 8D1?4-3
\textbf{Dạng 3. Rút gọn, chứng minh biểu thức không phụ thuộc vào biến}
\begin{enumerate}
\item Thực hiện các phép nhân đơn thức với đa thức, đa thức với đa thức.
\item Bỏ ngoặc và thu gọn toàn bộ các hạng tử chứa biến.
\item Nếu kết quả thu gọn là một hằng số (không còn chứa biến), ta kết luận giá trị của biểu thức không phụ thuộc vào giá trị của biến.
\end{enumerate}

%% DANG 8D1?4-4
\textbf{Dạng 4. Tìm $, bài toán thực tế (diện tích, thể tích)}
\begin{enumerate}
\item \textbf{Bài toán tìm $}:
\begin{itemize}
\item Khai triển phép nhân đa thức ở hai vế.
\item Thu gọn các hạng tử đồng dạng chứa ^2, x$. Thường các hạng tử bậc hai ^2$ sẽ triệt tiêu đưa về phương trình bậc nhất  + b = 0$.
\item Tìm $ và kết luận.
\end{itemize}
\item \textbf{Bài toán thực tế}:
\begin{itemize}
\item Biểu diễn các kích thước (dài, rộng, chiều cao) theo biến $.
\item Lập biểu thức diện tích hình chữ nhật ( = a \cdot b$) hoặc thể tích hình hộp chữ nhật ( = a \cdot b \cdot h$).
\item Khai triển và thu gọn đa thức.
\end{itemize}
\end{enumerate}
